\section{Théorème de Bézout et PPMC}

\begin{theoreme}[\textit{Théorème de Bézout}]%\label{coefbezout}
Soient $a;b \in \Z$. Alors il existe des entiers $\alpha$ et $\beta$ tels que
$$\pgdc(a;b)=\alpha\fois a + \beta\fois b$$
\end{theoreme}

\medskip
La preuve de ce théorème découle directement de l'algorithme d'Euclide.

\medskip
Les entiers $\alpha$ et $\beta$ s'appellent les \textbf{coefficients de Bézout} et s'obtiennent en \guillemets{remontant} l'algorithme d'Euclide.
Il faut noter que les entiers $\alpha$ et $\beta$ ne sont pas uniques.

\bigskip
\begin{exemple}
 Calculons les coefficients de Bézout pour $a=600$ et $b=124$.
 
\bigskip
{\small
$$
\begin{tikzpicture}
  \matrix (eucl) [matrix of math nodes, inner sep=0pt,column sep=4pt, row sep=4pt]{
      600 & = & 124 & \fois & 4 & + &    104     & \qquad & 4 & = &|[right]|
      \left[
        \begin{array}{l}
          600 \fois \textcolor{blue}{6} + 124 \fois (\textcolor{blue}{-29}) \\
          124 \fois (-5) + (600-124\fois 4) \fois 6                         \\
        \end{array}
      \right.
      \\
      124 & = & 104 & \fois & 1 & + &     20     &  & 4 & = &|[right]|
      \left[
        \begin{array}{l}
          124 \fois (-5) + 104 \fois 6\\
          104 - (124-104\fois 1) \fois 5 \\
        \end{array}
      \right.
      \\
      104 & = & 20  & \fois & 5 & + & |[red]|{4} &  & 4 & = & |[anchor=base west]|
      \left[
      \begin{array}{l}
        104 - 20 \fois 5
      \end{array}
      \right.
      \\
      20  & = &  4  & \fois & 5 & + &     0      & \qquad &   &   & \\
  };
  \draw[->, red, very thick] (eucl-1-8.west)..controls ([yshift=-1em]eucl-4-8.south west) and ([yshift=-1em]eucl-4-8.south east)..(eucl-1-8.east);
\end{tikzpicture}
$$} 

Les coefficient de Bézout sont donc $6$ et $-29$ et on a:
$$\pgdc(600;124)=4=\bm{6}\fois 600 + (\bm{-29})\fois 124$$
\end{exemple}


\smallskip
\begin{corolaire}\label{cor:diviseur_pgcd}
Si $d|a$ et $d|b$ alors $d | \pgdc(a,b)$.
\end{corolaire}

\smallskip
\begin{exemple} $4|16$ et $4|24$ donc $4$ divise $\pgdc(16,24)=8$.
\end{exemple}

\smallskip
\begin{demonstration}
Par le théorème de Bézout, il existe $\alpha,\,\beta\in\N$ tels que $\alpha\fois a + \beta\fois b=\pgdc(a,b)$.

\medskip
Comme $d|a$ et $d|b$ alors $d | (\alpha\fois a+\beta\fois b)=\pgdc(a,b)$ (combinaison linéaire).
\vspace{-3mm}\hfill $\blacksquare$
\end{demonstration}

\bigskip
\begin{corolaire}
Soient $a, b$ deux entiers. $a$ et $b$ sont premiers entre eux
\textbf{si et seulement si} il existe $u,v \in \Z$ tels que
$$ua+vb=1$$
\end{corolaire}

\begin{demonstration}
Le sens direct $(\Rightarrow)$ est une conséquence du corolaire \ref*{cor:diviseur_pgcd}.

\medskip
Pour le sens indirect $(\Leftarrow)$, supposons qu'il existe $u,v\in\Z$ tels que $ua+vb=1$.

\smallskip
Comme $\pgdc(a,b)|a$ alors $\pgdc(a,b)|au$. De même $\pgdc(a,b)|bv$.

\smallskip
Donc $\pgdc(a,b)|(au+bv)=1$. Donc $\pgdc(a,b)=1$.

\vspace{-3mm}\hfill $\blacksquare$
\end{demonstration}

\medskip
\begin{remarque}
Si on trouve deux entiers $u',v'$ tels que $u'a+v'b=d$, cela n'implique
\textbf{pas} que $d=\pgdc(a,b)$. On sait seulement alors que $\pgdc(a,b)|d$.

\smallskip
Par exemple si $a=12$, $b=8$ alors $12 \cdot 1 + 8 \cdot 3 = 36$ et $\pgdc(a,b)=4$.
\end{remarque}

\bigskip
\begin{corolaire}[\normalsize Lemme de Gauss]
\index{lemme!de Gauss}
Soient $a,b,c \in \Z$.

$$\textrm{Si} \ \ a | bc \ \  \textrm{\bf et} \ \  \pgdc(a,b)=1 \ \  \textrm{alors} \ \  a|c$$
\end{corolaire}

\medskip
\begin{exemple} si $4 | (7c)$, comme $4$ et $7$ sont premiers entre eux, alors $4|c$.
\end{exemple}

\begin{demonstration}
Comme $\pgdc(a,b)=1$, alors il existe $u,v\in \Z$ tels que
$ua+vb=1$.

\smallskip
On multiplie cette égalité par $c$ pour obtenir
$acu + bcv = c$.

\smallskip
Mais $a|acu$ et par hypothèse $a|bcv$ donc
$a$ divise $acu+bcv=c$.

\vspace{-3mm}\hfill $\blacksquare$
\end{demonstration}


\bigskip
\begin{definition}
Le \textbf{plus petit multiple commun} de deux nombres $a$ et $b$, noté $\ppmc(a,b)$,
est le plus petit entier positif divisible à la fois par $a$ et par $b$.
\end{definition}

\begin{exemple}$\ppmc(12,9)=36$.
\end{exemple}

Le $\pgdc$ et le ppmc sont liés par la formule suivante :

\medskip
\begin{proposition}
Si $a,b$ sont des entiers (non tous les deux nuls) alors,

\begin{center}
\fbox{$\pgdc(a,b) \cdot  \ppmc(a,b) = |ab|$}
\end{center}
\end{proposition}

\begin{demonstration}
Sans perte de généralité, on peut supposer $a>0$ et $b>0$.

\smallskip
Posons $d= \pgdc(a,b)$ et définissons $m= \frac{|ab|}{\pgdc(a,b)}$.

\bigskip
On écrit $a=da'$ et $b=db'$ avec $\pgdc(a',b')=1$. Donc 
$$ab=d^2 a'b'\eh{4}\textrm{et par suite}\eh{4}m=da'b'$$
Ainsi $m=ab'=a'b$ est un multiple de $a$ et de $b$.

\medskip
Il reste à montrer que c'est le plus petit multiple. Si $n$ est un autre multiple
de $a$ et de $b$ alors 

$$n= ka= \ell b$$

donc $$kda'=\ell d b' \eh{4}\textrm{et par conséquent}\eh{4} ka'=\ell b'$$

La dernière égalité implique que $a' | \ell b'$. Mais comme $\pgdc(a',b')=1$ , le lemmes de Gauss assure que $a'|\ell$ et donc $a'b | \ell b$.

\medskip
Comme $m=a'b$ et $\ell b = n$, nous avons que $m|n$.

\medskip
Nous avons montré que tout multiple commun de $a$ et $b$ est divisible par $m$, donc $m=\ppmc(a,b)$ et finalement 
$$\pgdc(a,b) \cdot  \ppmc(a,b) = |ab|$$

\vspace{-3mm}\hfill $\blacksquare$
 
\end{demonstration}

Voici un autre résultat concernant le ppcm qui se démontre en utilisant la décomposition en facteurs premiers (que nous étudierons à la section suivante):
\begin{proposition}
Si $a|c$ et $b|c$ alors $\ppmc(a,b)|c$.
\end{proposition}

\begin{remarque}
Il serait faux de penser que $ab|c$.

\smallskip
Par exemple $6|36$, $9|36$ mais $6\cdot 9$ ne divise pas $36$.
Par contre $\ppmc(6,9)=18$ divise bien $36$.
\end{remarque}

%---------------------------------------------------------------
%\subsection{Mini-exercices}

\medskip
\begin{exstheorie}
\exo
Calculer les coefficients de Bézout correspondants à :
\begin{enumerate}[label=\alph*), left=0mm,itemsep=2mm]
\item $\pgdc(7260,637)$.
\item $\pgdc(9945;3003)$. 
\end{enumerate}




\eexo





\exo
A l'aide d'un corollaire du théorème de Bézout, montrer que :
$$\pgdc(a,a+1)=1$$




\eexo





\exo
Trouver une solution aux équations suivantes : 
\begin{tasks}[after-item-skip=2mm,item-indent=5mm,column-sep=1mm,label-offset=1mm](1)
\task $873x+578y=1$
\task $315x+572y=5$
\task $6x+20y=3$
\end{tasks}



\eexo





\exo
Trouver tous les couples $(a,b)\in\Z^2$ vérifiant $\pgdc(a,b)=12$ et $\ppmc(a,b)=432$.
\end{exstheorie}


